% Cantor集前四步C_0..C_3：C_n由2^n个长3^{-n}的闭区间组成，总长度1,2/3,4/9,8/27；区间下方数字为该步选择的三进制位(0左段,2右段)。
\begin{tikzpicture}[x=8.2cm,y=1cm,
  note/.style={font=\FigureNoteFont\fontsize{7.5}{9}\selectfont,text=FigureMuted},
  dig/.style={font=\fontsize{6.5}{7}\selectfont,text=FigureMuted},
  lab/.style={font=\fontsize{8}{9}\selectfont}]
% recursive drawing via explicit lists
\newcommand{\seg}[3]{\draw[FigureBlue,line width=3pt] (#1,#3)--(#2,#3);}
\seg{0}{1}{0}
\node[lab,anchor=east] at (-0.03,0) {$C_0$};\node[note,anchor=west] at (1.03,0) {$1$};
\foreach \a/\d in {0/0,0.666667/2}{\seg{\a}{\a+0.333333}{-0.7}\node[dig] at (\a+0.166667,-0.98) {$\d$};}
\node[lab,anchor=east] at (-0.03,-0.7) {$C_1$};\node[note,anchor=west] at (1.03,-0.7) {$\tfrac23$};
\foreach \a/\d in {0/0,0.222222/2,0.666667/0,0.888889/2}{\seg{\a}{\a+0.111111}{-1.4}\node[dig] at (\a+0.055556,-1.68) {$\d$};}
\node[lab,anchor=east] at (-0.03,-1.4) {$C_2$};\node[note,anchor=west] at (1.03,-1.4) {$\tfrac49$};
\foreach \a/\d in {0/0,0.074074/2,0.222222/0,0.296296/2,0.666667/0,0.740741/2,0.888889/0,0.962963/2}{\seg{\a}{\a+0.037037}{-2.1}\node[dig] at (\a+0.018519,-2.38) {$\d$};}
\node[lab,anchor=east] at (-0.03,-2.1) {$C_3$};\node[note,anchor=west] at (1.03,-2.1) {$\tfrac{8}{27}$};
\foreach \x/\l in {0/0,0.333333/\tfrac13,0.666667/\tfrac23,1/1}{\draw[FigureGrid,dashed,line width=.5pt] (\x,0.15)--(\x,-2.5);\node[font=\fontsize{7}{8}\selectfont] at (\x,0.42) {$\l$};}
\node[note,anchor=west] at (1.03,0.42) {总长度};
\end{tikzpicture}
